\documentclass[a4paper]{article}

\usepackage{amsmath}
\usepackage{amsfonts}
\usepackage{amssymb}
\usepackage{graphicx}
\usepackage[spanish]{babel}
\usepackage[utf8]{inputenc}
\usepackage{hyperref}

\usepackage{authblk}

\title{Título del trabajo (en lo posible menos de 20 palabras)}

% Author names and affiliations
\author[1,2,*]{Charly García}
\author[2]{Pedro Aznar}
\author[1,+]{David Lebón}
\author[3,+]{Oscar Moro}

\affil[1]{Departamento de Matematicas, FCEN-UBA}
\affil[2]{Departamento de Fisica, FAMAF-UNC}
\affil[3]{Departamento XXX, Universidad YYY}

\affil[+]{Aclaracion sobre estos autores}
\affil[*]{\underline{direccion.de.correo\_del\_autor@servidor.com}}


\date{1/3/2018}
\setcounter{Maxaffil}{0}
\renewcommand\Affilfont{\itshape\small}

\begin{document}
\maketitle

\begin{abstract}

Acá va el texto del resumen. Podés usar los modificadores de \LaTeX 
para incluir \textbf{texto en negrita}, o \textit{en cursiva}, 
o \underline{texto subrayado}.

Mostrar que podemos poner fórmulas como:
\[ 
x^2+\phi = \int_0^{x} \frac{\theta(y)}{\pi} dy
\]

% Las lineas que empiezan con % son comentarios y no aparecen 
% impresas en el documento final


oindent
O hacer listados:

\begin{itemize}
  \item Item 1
  \item Item 2 que puede contener un \href{https://es.wikibooks.org/wiki/Manual_de_LaTeX}{link a un manual de Latex on-line}.
  \item etc
\end{itemize}

Ademas podes citar \cite{lamport94} documentos. Podes citar el Wolf \cite{Born1970} o
citar varias a la vez \cite{lamport94,Born1970}.

\end{abstract}



\renewcommand{\refname}{\small Referencias:}
\small
\begin{thebibliography}{9}

\bibitem{lamport94}
  Leslie Lamport, \textit{\LaTeX: a document preparation system},
  Addison Wesley, Massachusetts, 2nd edition, 1994.

\bibitem{Born1970}
  Born, M., \& Wolf, E. (1970). Principles of Optics. Optics \& Laser Technology. 
  University Press, Cambridge. 
  Retrieved from http://linkinghub.elsevier.com/retrieve/pii/S003039920000061X


\end{thebibliography}




\end{document}
